\section{Information categories and relation to earlier inverses}
\label{sec:comparison}
The principal theorem recovers the obstacle configuration and the
complete stationary launch density from a directional short-flight
germ. Its finite version regularizes this identity and prices the
actual forward collision bits. The fixed-known-disk minimax result
provides a separate benchmark. The appendices retain the pooled
occupation and homothetic constructions, with their distinct data and
hypotheses. We describe the classical ingredients and the additional
identification steps for each experiment.

\subsection{Directional boundary measures and the unknown density}
First variation of translated-set probabilities has established
geometric antecedents. Galerne~\cite{Galerne} relates directional
covariogram derivatives to directional variation and perimeter, and
Kiderlen and Rataj~\cite{KiderlenRataj} study first variations of
dilation volume and contact distributions of stationary random sets.
Lemma~\ref{lem:single-boundary-flux} derives the spatially weighted
entering-boundary measure for the present attempted-collision law,
including a weak finite-length error that uses only the total mass
of the unknown density.

The planar angular differential inverse is also classical. In the
sine-transform normalization, Louis, Riplinger, Spiess and
Spodarev~\cite[Section 4]{LRSS} give $f=(g+g'')/2$ for the even
angular datum. Our one-sided kernel has the normalization
\eqref{eq:single-cosine-distribution}. It leaves the sum of two
antipodal normal contributions. Angular inversion alone therefore
does not identify a nonsymmetric obstacle or an unknown launch law.
The spatial variable supplies the additional information: the positive
angular field consists of translated density copies. Their component
supports and masses identify the common footprint and density; their
Steiner points trace the obstacle boundaries. The minimum-area
selection in Theorem~\ref{thm:single-resolved-obstacle} extends the
exact inverse to configurations with only one sufficiently large
obstacle. The proof classifies the full geometric and probabilistic
fiber and determines the complete collision law from its first-order
germ.

The finite inverse uses positive regularization and places derivatives
on known test functions, as in classical approximate inversion.
The collision-specific quantitative steps are the weak flight bias,
density-uniform spatial quadrature, the averaged rare-event variance,
and separated-copy acquisition. The resulting bound is an explicit
sufficient cost for unknown-law reconstruction. The sharp common-disk
power of Section~\ref{sec:sharp-stationary} is stated on its own
observation and nuisance class.

\subsection{Unknown probes and global response rigidity}
Geometric inference with an unknown probe has a close precedent in
scanned-probe microscopy. Villarrubia~\cite{Villarrubia} expresses
probe imaging through morphological dilation and develops blind
reconstruction from image features. Its consistent-probe envelope
is an important comparison for an unknown Minkowski summand.
In the retained homothetic construction, the extra collision flux
identifies footprint width, and two translated settings give exact
support-envelope cancellation. The resulting conclusion identifies the entire
geometric ambiguity, without asymmetric species, correspondence or a
periodic presentation.

The stopped Poisson representation and discrete obstacle iteration
use classical potential theory \cite{LawlerLimic}. Their role in
Theorem~\ref{thm:global-occupation} is to recover the positive set
without a supplied component or killing domain. Uniformly bounded
separated components give the stopping budget, and support
cancellation yields the exact equality of period groups. This
whole-field statement differs from the finite protected-patch
period tests; its finite-field locality bound explicitly records
how the observation region grows.

\subsection{Boundary estimation and informative rare responses}
Active boundary estimation is an essential comparison for the finite
theorem, not merely an analogy with passive billiard spectra.
Castro and Nowak~\cite{CN} study adaptive binary classification under
boundary-fragment smoothness and noise assumptions, with matching
orders for their excess-risk criterion. Their input selects a feature
point whose response is a class label. This directly available label
is not our collision bit: a solid start and a free miss have the same
output here. Neither their observation model nor their loss identifies
our finite periodic table. We use an independent physical packing and
binary-tree proof rather than importing a classification lower bound.

Locatelli, Carpentier and Kpotufe~\cite{LCK} give an adaptive strategy
for smooth decision boundaries without prior knowledge of smoothness
and noise parameters. Their membership-query setting and local
line-search/interpolation construction make the closest algorithmic
comparison. In contrast, our smoothness is known, our error is a $C^2$
geometry norm, and the boundary label must first be derived from
pooled collision data. We claim neither a new general
bisection principle nor an advance on adaptation to unknown smoothness.
The reductions in Propositions~\ref{prop:query} and
\ref{prop:rare-query} permit this active-estimation strategy in the
physical collision experiment, at their respective costs.

Yan, Chaudhuri and Javidi~\cite{YCJ} use informative abstention
responses in active learning. Under their assumptions the rarity of
an informative response can yield an inverse-probability cost. Their
observation includes an abstention symbol. Our sensor has only its
original collision bit. The additional argument here is that translated
pooled commands and a conjunction over outward compass candidates
realize the requisite one-sided event for each boundary query. Small
occupation alone would not justify such an improvement: the individual
means in a signed reciprocal difference need not be small. The exterior
zero and interior lower bound are therefore proved for the commands
that are actually sampled.

Weighted cap mass is also classical. Brunel~\cite{BrunelCaps} gives
the exponent $\gamma+(d+1)/2$ for rolling convex supports and
boundary-decaying densities, and translates cap assumptions into
Hausdorff rates for random convex hulls of observed spatial points.
In the plane the geometric exponent is $\gamma+3/2$. The present
upper bound must realize such information from hidden-launch bits.
The sharp lower bound additionally controls the actual swept-minus-solid
response: its effective-boundary cancellation is uniform for every
short direction and for lengths approaching zero. A cap probability
alone would not prove that all such commands have the required
Hellinger contraction.

\subsection{Localization and control precision}

The earlier scalar inverse reconstructed a uniformly accurate occupation
on a full two-dimensional grid. Its constructive upper bound was
$C\nu^{-3}\log(C/(\nu\delta))$ attempts on the same bounded
$C^{6,\beta}$ class, using $\sigma\asymp\nu^{3/2}$ and compensated
support smoothing. That theorem remains in the exact repository archive of A2 v31. The retained adaptive localized
procedure uses $O(h^{-1})$ angular lines, a fixed-depth dependency
stencil for each queried point, and the full $C^{6,\beta}$ regularity.
It has a different adaptive command design, not a sharper analysis of
the old uniform design. We have not proved a matching lower bound
for that old design or for uniform spatial preparations.

The resource improvement is accompanied by the explicit localization
and calibration cost in \eqref{eq:precision}. In that localized experiment the exact data are a pair
of functionals on controlled spatial inputs. The lower bound concerns
only their finite binary realizations; it does not assert a low-information
passive invariant. The angle/time/position tolerances shrink with
localization and are not learned from the collision outcomes. Apparatus
motion, precision in representing input locations, and arithmetic
cost are distinct resources. Goldberg and Kwek~\cite{GK} treat
query-coordinate precision as an explicit resource in learning convex
polytopes by exact membership queries. The digital bounds in
Theorems~\ref{thm:unk-scale-finite} and \ref{thm:registration-finite}
follow the same accounting principle
for a different class and observation: it includes the additional
integral-measurement centers as well as the fine angular boundary
queries. It does not convert a bit-description bound into a physical
metrology or travel bound.

For completeness the exact v31 manuscript is preserved unchanged
inside the repository archive of the reviewed v33 source. It retains the arbitrary centered-law and general
nontrivial bounded-law inverses, the different-zero-set comparison,
the full fixed-aperture proof, rational interval enclosures, the
uniform-grid construction, and all earlier moment, marked-law and
count-fiber results in their original nested volumes. The more general
randomized law is an exact functional theorem; the finite local query
proved here uses the four-atom compass law. No continuous transition
law is discretized without an error budget. Rational intervals remain
conditional on forcing and survival bounds, not certificates of the
apparatus or physical prior.

The finite periodic conclusion also retains the known positive patch
margin and the bounded periodic presentation. Exact rational period
relations are distinct from approximate Euclidean generators. Removing
the known margin gives only pointwise eventual recognition, as described
in Section~\ref{sec:periods}. The exact full-field period criterion in
Theorem~\ref{thm:global-response-rigidity} has no periodicity prior,
and Corollary~\ref{cor:registration-finite-cloud} treats a fully
captured finite configuration without a period margin. The direction-resolved forward germ of
Section~\ref{sec:single-law-rigidity} has prescribed nominal positions
and directions. Uniformly prepared count germs, analytic network
records, marked length spectra and passive trajectory laws are
different observations. The hypotheses of each inverse specify
which of these data are available.

\subsection{Stationary noise versus adversarial implementation}
The retained fixed-footprint construction belongs also to support estimation with
contaminated data. Brunel, Klusowski and Yang~\cite{BKY} observe
continuous random vectors drawn from an unknown convex body and then
contaminated by additive Gaussian noise. Their geometric estimator
avoids Fourier inversion and is analyzed in Hausdorff loss. Here no
random vector or direct inside/outside response is observed: the data
are attempted collision bits at controlled centers. The fixed compact
noise support, its boundary-mass lower bound and the reciprocal reduction
are essential to our polynomial upper bound. The Gaussian observation
and its tail-based analysis are not covered by this theorem, and its
rates are not being compared as rates in a common experiment.

Support addition, inner rolling disks and the Brunn--Minkowski inequality
are classical \cite{Schneider}. Their role in the retained occupation construction is to remove an
unknown stationary density without learning that density. The positivity set of
the recovered blur, not its detailed values, identifies the footprint
sum. The retained signed high-accuracy query uses the Green bound for exit from a fixed
rectangle; its constants can be large and its finite-depth work grows
logarithmically. The area normalization uses a finite adjoint of the
same killed operator. The direct collision-strip normalization instead
uses Cavalieri's principle and width additivity; only one integrated
pooled forward mean is required. We do not claim these convex-geometric
or random-walk
principles as new abstract theorems.

The localized and stationary upper bounds cannot be ranked by their
attempted-bit exponents alone. The localized model attains the smaller
sample power but requires shrinking spread and worst-case tolerance.
The stationary model keeps its random spread fixed and allows the
density to be unknown. Its uniform-disk minimax power is sharp by
Corollary~\ref{cor:stationary-minimax}; the comparison uses the same
physical class, confidence and expected-cost criterion, with a
stronger command class allowed for the lower bound. For the rare-query
upper, $2t+\max_i\diam A_i<d_0$ is displayed explicitly; the older
signed-query upper needs only $t+\max_i\diam A_i<d_0$.
The retained homothetic geometric inverse reconstructs the centered footprint,
ratios and relative setting translations without a supplied common
origin. Exact homothety, labelled settings, per-setting stationarity
and nominal positioning remain part of its observation model.
Quantitative boundary-mass, geometry, scale-gap and coarse offset
bounds enter its uniform finite version. The common-disk converse
does not extend automatically to every stationary density or to
vanishing launch spread, and digital command counts do not assign
manufacturing or metrology costs.
